% !TeX root = ../papers/nc-p3.tex
\section{階数二曲線上の四次曲面と八次判別式の二重接触}
\label{sec:cofactor-quartic-contact}
定理\ref{thm:cofactor-no-integral-cubic}の仮定を用い、一般被約な余因子曲線の整な台を $C$ とする。
\begin{lemma}\label{lem:cofactor-linear-gluing}
整曲面 $F\subset\mathbb P^3$ の次数を三以上とする。一般平面 $H$ の切断上で
有理写像 $\phi:F\dashrightarrow\mathbb P^n$ が $H$ 上の一次写像の制限と一致し、
$\phi$ が非定数なら、$\phi$ は周囲の一つの一次写像の制限である。
\end{lemma}
\begin{proof}
二つの一般平面の交線は $F$ を相異なる三点以上で切る。
そのうち二点の像は一般に異なる（一般二点の割線は一般直線である）。
各平面上の一次写像は交線で階数二であり、この三点の値によって
交線全体で同じ射影写像となる。従って一方の定数倍を調整すれば両線形写像は交線上で一致し、
二平面から周囲の一つの線形写像 $A$ を構成できる。
第三の一般平面で、先の二平面との各交線を考えれば、その上の一次写像は両交線で $A$ と定数倍で一致する。
三平面の交点では $A$ は非零なので二定数は等しく、第三平面でも $A$ が使える。
基点がある場合も、一般に選んだ交線上で二点の像が異なればそこに基点はない。
\end{proof}

\begin{theorem}\label{thm:cofactor-quartic-contact}
$C$ を含む四次曲面が存在するならただ一つで、整である。このとき
\[
        \deg C=16,\qquad [F\cap V(\det N)]=2[C] .
\]
この四次曲面は行列式の因子にならない。一般平面切断 $E=F\cap H$ と $\Gamma=C\cap H$ に対し
$\mathcal O_E(4)(-\Gamma)$ は非自明な二捩れである。
従って八次行列式には $C$ に沿う重複度二の既約因子がただ一つあり、その因子自体の次数は五以上で、
行列式の因数分解における指数は一である。
\end{theorem}
\begin{proof}[四次因子の排除]
含む四次式 $F$ は前節より既約であり、その特異軌跡は $C$ を含まない（偏導は三次）。
$J=\mathcal I_{C/F}$ として、前節と同じ双対化により
$H^0(\mathcal Hom_F(J,\mathcal O_F(-3)))=H^0(\omega_C(-3))$ は二次元以上。
独立な二切断 $u_1,u_2$ を取り、整な一般平面四次 $E$ 上で
$M_E=\mathcal O_E(\Gamma-3H)$ に制限すると独立である。ここで $H$ は超平面類も表す。
$M_E$ の次数は $d-12$ で、二以上の独立切断を持つ線束の次数は三以上である。
実際次数二の有効Cartier因子は長さ二で、標準束 $\omega_E=\mathcal O_E(1)$ がその長さを分離するため、
Riemann--Rochで切断は一つである。低い次数は滑らかな点を足して帰着する。よって $d\ge15$。

$F\mid\det N$ と仮定する。一般平面四次の特異点は階数三なので前節と同じく核射が $E$ 上で定義され、
その線束 $L$ は $2\deg L\le24-d\le9$ を満たす。核像が一般 $E$ で直線に入れば、
二一般平面を比較する前節の議論により $F$ 全体で一固定直線に入る。定数核も同様に除ける。
直線の場合は $Nv\wedge Nw$ が既約四次 $F$ と定数外積の積となり、
一般可逆性と一列の四成分の独立性に反する。従って $h^0(L)\ge3$。
次数四以下のこのような線束は $\mathcal O_E(1)$ だけである。
実際次数四ではRiemann--Rochと次数零の双対線束の切断を使い、次数三以下は標準束に埋め込めば
一点の消滅で切断が二次元以下になる。従って核点は各一般 $E$ で平面上の一次写像で書ける。
補題から $F$ 上で $Ax$ が核点となり、三次列 $N(x)Ax$ は四次曲面上零、従って恒等零である。
これは一般可逆性に反する。以上でどの四次曲面も因子になれない。

一般 $E$ で行列式切断は非零かつ $\Gamma$ で二重なので $2d\le32$、従って $d=15,16$。
また行列式の $C$ における位数はちょうど二である。二行を消した残余対称二行の三成分は、
一般被約性から二次元正則局所環の極大イデアルを生成する。それらの一次部の行列の像の射影は直線で、
二行対称の行列式円錐の射影二次曲線に直線はないから初項の二次式は非零である。
$C$ を含む各既約因子は次数五以上であり、二つでもその一つの平方でも次数八を越える。
従ってただ一因子が指数一で $C$ に沿って二重である。この結論は $C$ を含む四次式が存在しなくても成立する。
\end{proof}

\begin{proof}[二つの四次式と残余直線]
独立な二つの四次式 $f,g$ が $C$ を含むと仮定する。一般に取り直せば、鉛筆の一般曲面
$F=V(f)$ の一般平面切断 $E$ は整である。$d=16$ のとき完全交叉は次数から $C$ に一致し、
定理\ref{thm:cofactor-no-integral-ci}に反する。$d=15$ なら完全交叉での残余は純次数一の直線 $T$。
完全交叉と $C$ がともにCohen--Macaulayであることによる。従って一般点では $C,T$ の両方で
$f,g$ の微分は独立、一般平面上で $g|_E$ の零因子は $\Gamma+p$ ($p=T\cap H$) である。
$\det N/g^2$ は $E$ 上で $2p$ 以内の極しか持たず、
上で示した次数二の切断数より定数である。一般平面上で定数なので $F$ 上でも定数で、
$\det N= c g^2+f h$、$h$ は四次である。$C$ の一般点で左辺と $g^2$ は二重に消え、
$f$ は単純なので $h$ は $C$ を含む。
このとき $I_C(4)$ の空間は $f,g$ の包である。実際 $h|_E/g|_E$ は高々単純な一点の極しか持たず、
一般 $E$ で定数となり、$F$ でも定数である。
従って行列式は $f,g$ の定数係数二次式で、分解して $C$ を含む四次因子を持つ。前段と矛盾する。
これで四次曲面の一意性を得た。
\end{proof}

\begin{proof}[十五次と自明な接触束の排除]
$d=15$ なら $M_E$ は次数三で二切断を持つ。双対性から
$M_E=\mathcal O_E(1)(-p_E)$ で、$p_E$ は滑らかな一点である
（長さ一の有効Cartier因子の点は正則である）。
その二切断はその点を通る直線の完全系で、基点を除いた束としては大域生成される。
$d=16$ なら $\det N|_E$ の零因子はちょうど $2\Gamma$ なので
$\mathcal O_E(4)(-\Gamma)$ は二捩れである。それが一般 $E$ で自明な場合を合わせて考える。
そのとき $M_E=\mathcal O_E(1)$ なので二切断の比も一次式の比である。
いずれも補題より $u_1/u_2=l_1/l_2$ と書け、$l_1,l_2$ は周囲の独立な一次式である。
$u=u_1/l_1=u_2/l_2$ は $F$ 上の次数 $-4$ の有理切断である。
$E$ では一次式二つの唯一あり得る共通零点は $p=H\cap T$、$T=V(l_1,l_2)$ である。
$d=15$ なら上の完全系の次数からこの共通零点はちょうど $p_E$ であり $T\subset F$。
$d=16$ の場合にも $T\subset F$ はあり得るが、そのとき $u_1,u_2$ にも同じ共通零がある。
実際 $M_E=\mathcal O_E(1)$ の同型の下で二組の一次式は射影的に一致し、
平面四次上での二次の交代積の消滅から平面上でも定数倍で一致する。

従って $u|_E$ は次数 $-4$ の有理切断として $\Gamma$ で単純な極を持ち、
$d=15$ ではさらに $p_E$ で単純な極を持つだけで、零はない。
より正確に、$\mathcal O_E(\Gamma)$ を用いて $u_i$ を正則切断としてから一次式で割った記述なので、
平面切断の特異点でも残りの係数は単元である。
この情報は $1/u$ が $F$ の全ての余次元一点で正則で、$C$ に沿って消えることを与える。
実際 $C$ と十五次の場合の $T$ では一般横断平面のその点が滑らかなので離散付値を読める。
他の曲線の一般点では $J=\mathcal O_F$ で、$T$ を除けばどちらかの $l_i$ が単元なので $u$ 自体が正則、
一般平面への制限が単元だから逆も正則である。十六次で $T\subset F$ の場合も、
一般平面で $\mathcal O_E(1)$ の二切断が一次式定数倍であるためこの結論が要る。
次段の局所的な議論でその場合を含める。

一般閉点 $x$ で $1/u$ が属するか調べる階数一の有理切断の分母イデアルを取る。
分数 $1/u=A/B$ を局所正則関数で表し、$B$ は非零とする。
$A\bmod B$ が曲線成分の一般点で非零なら、その生成する部分加群は一般横断平面を法としても
$\mathcal O_F/(B)$ の中で非零である。実際一般平面式は同部分加群と商の随伴素点を避け、
テンソル後の単射を保ち、交点で中山の補題が使える。
従って一般平面上で分数が正則であることは元の余次元一での正則性を保証する。
$F$ は超曲面で $S_2$ なので $1/u$ は $\mathcal O_F(4)$ の大域切断に延長し、
四次式へ持ち上がる。これは $F$ と独立で $C$ を含み、一意性に反する。
ゆえに十五次はなく、十六次の二捩れは一般平面切断で非自明である。次数と二重消滅から
記載した交叉サイクルの等号も従う。
\end{proof}
